% Section 3 -- Model (outline: paragraphs and what each must say; no prose yet)
%
% Written from the code of the R0 model, not from the pd5 draft:
%   spec        model-codes/codes/epc-lifecycle-deep-learning/epc_revised/spec.py
%               (build_model("pnr1_revised_2026"))
%   economics   estimation_epc/model_funcs_estimation_epc.py, epc_revised/economics.py,
%               epc_backward/physics.py, epc_backward/continuation.py
%   values      replication/data/model/r0/parameters.json,
%               replication/output/tables/tab-fixed-parameters.csv, tab-estimates.csv
% The pd5 model text (sections-pd5/03-model.tex) describes an older model; do
% not copy from it without checking against the code.
\section{Model}\label{sec:model}

\outpara{Section opening (one paragraph)}
\begin{itemize}
  \item What the model is: a partial-equilibrium lifecycle model of
    households who choose consumption, saving, renting or owning, the size
    of the home, mortgage debt, and the energy efficiency of the home they
    own.
  \item Efficiency is binary and matches the data: an efficient home is in
    bands A--C, an inefficient home in bands D--G (Section~\ref{sec:data}).
  \item The two explanations of Section~2 are two forces in the model:
    (i) the \emph{payoff horizon} -- how many years an owner expects to
    benefit from an efficient home; (ii) \emph{acquisition opportunities}
    -- efficiency is cheaper to obtain when buying a home than by renovating
    the home one lives in, and older owners buy less often.
  \item No preference for efficiency that varies with age: any age gradient
    comes from these forces.
  \item Road map: environment (3.1), household problem (3.2), the forces
    behind the age gradient (3.3). Parameter values in
    Section~\ref{sec:estimation}; full details in Appendix~B.
\end{itemize}

% ----------------------------------------------------------------------
\subsection{Environment}

\outpara{Demographics}
\begin{itemize}
  \item One period is one year. Households enter at age 23, work until 64,
    and retire at 65.
  \item Survival probability $\phi_j$ by age (ONS life table, 2015--17).
  \item Ages 82 and over form one group with a constant annual survival
    probability (0.880): no one is forced to die at a fixed age. Footnote:
    this avoids a terminal age at which owners would sell or stop renovating
    for mechanical reasons.
  \item Equivalence scale $e_j$ for household size (WAS).
  \item \open{Symbol clash: the table of estimates uses $\phi$ for the
    comfort value of efficiency and the fixed-parameter table uses $\phi_j$
    for survival. Proposal: keep $\phi_j$ for survival and rename comfort
    (e.g. $\eta$) in \codepath{tab_estimates.R} (exhibit code only, no
    change to estimation).}
\end{itemize}

\outpara{Preferences}
\begin{itemize}
  \item Equation: expected lifetime utility = sum of discounted period
    utility, plus warm-glow bequest at death, plus taste shocks on the
    discrete choice.
  \item Equation: period utility
    $u_j(c,s)=e_j\,\frac{(c^{\alpha}s^{1-\alpha})^{1-\sigma}-1}{1-\sigma}$,
    $\alpha=0.75$, $\sigma=2$.
  \item Equation: bequest
    $g(q)=\upsilon\frac{(q+\bar q)^{1-\sigma}-1}{1-\sigma}$;
    the estate $q$ is liquid wealth plus the market value of the home
    (at its efficiency at death) minus the mortgage. No selling cost.
    The efficiency premium therefore reaches heirs only if the home is
    still efficient.
  \item Taste shocks: Type-I extreme value, scale $\sigma_\varepsilon=0.10$,
    one per discrete alternative, iid over time. Gives logit choice
    probabilities.
\end{itemize}

\outpara{Income}
\begin{itemize}
  \item Working life: $y_j=\kappa_j z_j \nu_j$; permanent component a
    random walk, transitory iid; both mean-one lognormal
    ($\sigma_p=0.096$, $\sigma_\nu=0.206$; Understanding Society).
  \item Retirement: pension $=\rho_{\mathrm{ret}}\kappa_{64}z_{64}$, with a
    floor and a cap; $\rho_{\mathrm{ret}}$ is estimated.
  \item Net income $\tau_0 y^{1-\tau_1}$ (progressive taxes).
  \item One sentence on units: a unit of the model is mean working-age
    income (\pounds49{,}914); prices and sizes are in these units.
  \item Keep the income equations compact (one display, three lines).
\end{itemize}

\outpara{Housing and energy efficiency}
\begin{itemize}
  \item A home has a size $h$ (a quantity of housing) and an efficiency
    state: efficient ($\R$, A--C) or inefficient ($\NR$, D--G).
  \item Prices are constant: $p^{\NR}=1$ (numeraire) and
    $p^{\R}=(1+\gamma_p)p^{\NR}$, $\gamma_p=0.03$ (price premium of an
    efficient home).
  \item Equation: owner housing services
    $s^o(h,\R)=\left[\omega_0+\eta+\omega_h\log(1+h/h_\omega)\right]h$ and
    $s^o(h,\NR)=\left[\omega_0+\omega_h\log(1+h/h_\omega)\right]h$. The term
    $\eta$ (estimated) is the comfort value of an efficient home: a warmer
    home gives more housing services per unit of size.
  \item Recurring costs of an owner: maintenance $\delta_h p^q h$ and the
    energy bill $e^q h$, with $e^{\R}<e^{\NR}$ (EPC register and DESNZ
    prices). Efficient homes cost less to heat.
  \item Transaction costs: buyer pays $\zeta_b$, seller pays $\zeta_s$ of
    the value.
  \item Owned size between a minimum $h^{\min}$ (estimated; it governs who
    can afford to own) and a maximum.
  \item Summary sentence for the reader: efficiency pays in three ways --
    comfort, lower bills, and a higher value at sale or death.
  \item Interpretation: the housing stock in the model stands for homes
    built by 1990, the stock in which the Section~2 facts are measured.
    \open{confirm: the efficiency targets are for homes built by 1990 and
    the model has no dwelling vintage}
\end{itemize}

\outpara{Loss of efficiency}
\begin{itemize}
  \item An efficient owned home stays efficient next year with probability
    $\varrho$ and otherwise becomes inefficient. This applies to every
    efficient owned home, including one just bought or just renovated.
  \item Motivation: Fact~2 (downgrades in the register). Interpretation:
    wear of heating systems and insulation, and rising standards. One
    sentence.
  \item Inefficient homes become efficient only by renovation: upgrades are
    choices, downgrades are shocks.
  \item $\varrho$ is estimated; it sets the expected life of an efficiency
    investment, $1/(1-\varrho)$ years (about six years at the estimate).
\end{itemize}

\outpara{Three ways to obtain an efficient home}
\begin{itemize}
  \item (a) Buy a home that is already efficient: pay the price premium
    $\gamma_p$ (plus transaction costs on it).
  \item (b) Buy an inefficient home and renovate it in the year of
    purchase.
  \item (c) Renovate the home one already owns.
  \item Equation: renovation cost $C(h)=\xi_f+\xi\,p^{\NR}h$, the same in
    (b) and (c); $\xi_f$ fixed (estimated; search, disruption), $\xi=0.08$
    proportional (materials and work).
  \item Key point for the reader: a buyer can obtain efficiency by paying
    the 3 percent premium, while a stayer must pay the renovation cost
    (8 percent of value plus the fixed cost). A purchase is therefore a
    cheaper time to obtain efficiency. This is the ``renovation window'' of
    the title, and it works through the market for efficient homes, not
    through a discount on renovation at purchase.
  \item One sentence: there is no discount on the renovation cost at
    purchase ($\theta_{\mathrm{move}}=1$), so (b) costs more than (a) and is
    chosen rarely; it differs from (a) in collateral (the loan is set on
    the inefficient value) and in its taste shock.
  \item \open{Decision: keep (b) in the main text as one sentence, or
    describe only (a) and (c) and put (b) in Appendix~B. It is part of the
    estimated model either way.}
  \item \open{The Section~2 bridge says ``if renovation happens mostly
    around a purchase''; in the model, efficiency is mostly \emph{acquired}
    at purchase by buying an efficient home. Align the wording.}
\end{itemize}

\outpara{Renting}
\begin{itemize}
  \item Renters choose the size of the home each year; there is no
    transaction cost for renting.
  \item Rent per unit of size is fixed (gross yield of 5 percent).
  \item Renters do not choose efficiency: they rent from a stock in which a
    share $\pi_R=0.175$ is efficient (EHS private rentals). They pay the
    average energy bill and get the average comfort.
  \item Reason in one sentence: whether a rental is efficient is the
    landlord's decision, which is outside the model; this keeps the
    renovation decision with owners, where the facts are measured.
\end{itemize}

\outpara{Saving and mortgages}
\begin{itemize}
  \item Liquid saving at $r_b$; no unsecured borrowing ($b\ge 0$).
  \item Repayment mortgages at $r_m$; equation for the annual payment
    $\pi(m,n)$. A new loan has a 25-year term. Owners may repay faster
    than scheduled.
  \item Equation: at purchase,
    $m'\le\min\{\lambda_L p^q h',\ \text{loan whose payment}\le\lambda_D y^{P}\}$,
    with $\lambda_L=0.90$, $\lambda_D=0.33$; $y^{P}$ is gross income
    without the transitory shock.
  \item Movers who have a mortgage keep its remaining term, and can borrow
    at most the scheduled balance plus any increase in home value (plus
    the renovation cost if they renovate at purchase): no cash-out through
    moving. One sentence why: it rules out moves made only to refinance.
  \item An owner who renovates in place can borrow the cost, within the
    same limits, with the home valued as inefficient.
  \item Footnote: the premium of an efficient home cannot be borrowed
    against before the work is done.
  \item Details (term rules, edge cases) in Appendix~B.
\end{itemize}

\outpara{Initial conditions (short)}
\begin{itemize}
  \item Households enter at 23 with income, liquid wealth, tenure, home
    value and mortgage drawn from WAS heads aged 20--29.
  \item Initial owners are in an efficient home with probability 0.218
    (EHS, young owners).
\end{itemize}

% ----------------------------------------------------------------------
\subsection{The household problem}

\outpara{Timing within a year}
\begin{itemize}
  \item Start of year: income realised; cash on hand
    $x=(1+r_b)b+y^{\mathrm{net}}$.
  \item The household chooses a discrete alternative and, within it,
    consumption, saving, and (if buying) size and mortgage, or (if staying)
    the mortgage payment.
  \item End of year: the household dies with probability $1-\phi_j$ and
    leaves its estate; survivors draw next year's income, and an efficient
    home loses efficiency with probability $1-\varrho$.
\end{itemize}

\outpara{State and choices}
\begin{itemize}
  \item State: liquid wealth, mortgage balance and remaining term,
    permanent and transitory income, tenure and efficiency of the home
    (renter, inefficient owner, efficient owner), size of the owned home,
    and age.
  \item New table (choices): discrete alternatives by starting tenure.
    Renter (4): rent; buy inefficient; buy efficient; buy inefficient and
    renovate. Efficient owner (5): stay; sell and rent; sell and buy each
    of the three purchase options. Inefficient owner (6): the same plus
    renovate in place. 15 alternatives in all. Second column: the
    continuous choices in each.
\end{itemize}

\outpara{Bellman equations (two displays)}
\begin{itemize}
  \item Value of alternative $d$: best continuous choice of
    $u_j(c,s)+(1-\phi_j)g(q)+\beta\phi_j\,\E[V_{j+1}]$ subject to the
    budget and borrowing limits of $d$.
  \item Value before taste shocks: the expected maximum over alternatives,
    $\sigma_\varepsilon\log\sum_d\exp(Q^d_j/\sigma_\varepsilon)$, and the
    logit choice probabilities.
  \item Words after each equation: what each term is. No post-decision
    value function here; that belongs to the solution method
    (Section~\ref{sec:estimation} and Appendix~C).
\end{itemize}

% ----------------------------------------------------------------------
\subsection{What drives the age gradient in the model}

\outpara{The payoff from an efficient home}
\begin{itemize}
  \item Per year an owner gains comfort and saves
    $(e^{\NR}-e^{\R})h$ on bills; at sale or death the home is worth
    $\gamma_p p^{\NR} h$ more if it is still efficient.
  \item These gains last only while the owner lives, stays, and the home
    stays efficient. Define the payoff horizon (expected years of benefit):
    survival $\times$ retention $\times$ staying. This is the object
    reported in Section~\ref{sec:results} (Table on payoff horizons).
  \item Older owners have a shorter horizon because survival falls with
    age. A worked example in units of pounds for a typical home (bill
    saving per year against the cost of renovation) makes the trade-off
    concrete. \open{compute the example from R0 values: typical owned
    size, bill saving $\approx$ \pounds540 a year at $h=5$, renovation cost
    $\approx$ \pounds28{,}000 at $h=5$}
\end{itemize}

\outpara{Acquisition opportunities}
\begin{itemize}
  \item Efficiency is cheaper to obtain at purchase (premium) than in place
    (renovation cost). How often an owner buys therefore matters.
  \item Moving is costly (transaction costs, borrowing limits), and older
    owners have less reason to move, so they buy less often and hold the
    home they bought long ago -- whose efficiency has had time to decay.
  \item This is the model counterpart of Fact~3.
\end{itemize}

\outpara{Other forces}
\begin{itemize}
  \item Bequests: a stronger bequest motive raises the value of the
    efficiency premium for older owners (it passes to heirs) and makes
    them hold more housing wealth.
  \item Credit: borrowing limits bind for the young; the young buy more
    often but have less cash for the premium or a renovation.
  \item Comfort: the only utility gain of efficiency; it is the same at
    every age.
\end{itemize}

\outpara{What the model leaves out (short)}
\begin{itemize}
  \item Prices, rents and interest rates are fixed (partial equilibrium);
    no aggregate shocks.
  \item No differences between cohorts: every cohort faces the same
    environment. This is what lets the model separate age from cohort
    (Section~2 bridge).
  \item Landlords' decisions; the size of renovations (efficiency is
    binary); energy price shocks.
\end{itemize}

\outpara{Appendix~B (model details; not in the main text)}
\begin{itemize}
  \item Cash flows of each alternative; sale proceeds.
  \item Mortgage terms: porting, the borrowing cap of movers, the
    renovation loan, payment in retirement.
  \item The 82+ group and the simulation horizon.
  \item Initial conditions.
  \item Timing conventions (e.g. a household that renovates in place pays
    the inefficient energy bill in that year).
\end{itemize}
